\documentclass[12pt]{article}

\usepackage[margin=1in]{geometry}
\usepackage{amsmath, amssymb}
\usepackage[hidelinks]{hyperref}

\title{Feature Learning and MLP Manifold Muon}
\author{}
\date{\today}

\newcommand{\R}{\mathbb{R}}
\newcommand{\rms}{\operatorname{RMS}}
\newcommand{\Stiefel}{\operatorname{Stiefel}}
\newcommand{\Retr}{\operatorname{Retr}}
\newcommand{\msign}{\operatorname{msign}}
\newcommand{\fanin}{d_{\mathrm{in}}}
\newcommand{\fanout}{d_{\mathrm{out}}}

\begin{document}
\maketitle

\section{Purpose}

This note records the current MLP-only interpretation of how Manifold Muon
relates to hyperparameter transfer and the spectral feature-learning line of
work. It focuses on three claims:
\begin{enumerate}
    \item the feature-learning requirement from
    \href{https://arxiv.org/abs/2310.17813}{A Spectral Condition for Feature Learning};
    \item why Manifold Muon with no outer scale factor can nevertheless show
    strong hyperparameter transfer;
    \item how the $\sqrt{\fanin/\fanout}$ RMS-to-RMS factor fits the spectral
    scaling story.
\end{enumerate}

Throughout, consider an MLP hidden-layer matrix
\[
    W \in \R^{\fanout \times \fanin}.
\]
The discussion is for hidden matrix layers. Embeddings, biases, scalar
parameters, and final classifier heads require separate conventions.

\section{Feature-Learning Requirement}

For a linear layer, the first-order change in preactivation is
\[
    \Delta y = \Delta W x.
\]
Thus the natural object is not only the coordinatewise size of $\Delta W$, but
the size of $\Delta W$ as an operator acting on activations.

The spectral feature-learning condition asks the weight and the update to have
spectral norm on the order
\[
    \|W\|_2 = \Theta\!\left(\sqrt{\frac{\fanout}{\fanin}}\right),
    \qquad
    \|\Delta W\|_2 =
    \Theta\!\left(\sqrt{\frac{\fanout}{\fanin}}\right).
\]
This makes the induced change in features remain nontrivial and non-exploding
as width changes.

\subsection{RMS-to-RMS Form}

If vectors are measured in RMS norm,
\[
    \rms(x) = \frac{\|x\|_2}{\sqrt{d}},
\]
then the induced RMS-to-RMS operator norm of a matrix is
\[
    \|M\|_{\rms \to \rms}
    :=
    \max_{x \ne 0}
    \frac{\rms(Mx)}{\rms(x)}
    =
    \sqrt{\frac{\fanin}{\fanout}}\|M\|_2.
\]
So the same feature-learning condition can be written as
\[
    \|W\|_{\rms \to \rms} = \Theta(1),
    \qquad
    \|\Delta W\|_{\rms \to \rms} = \Theta(1).
\]

This is the source of the $\sqrt{\fanin/\fanout}$ factor. It is a
norm-conversion factor. If a tangent direction $A$ has $\|A\|_2 \le 1$ and we
want the ambient update to have RMS-to-RMS size $\eta$, then the update is
\[
    \Delta W
    =
    -\eta
    \sqrt{\frac{\fanout}{\fanin}}
    A.
\]
The norm conversion contains $\sqrt{\fanin/\fanout}$, while the tangent
multiplier contains its reciprocal $\sqrt{\fanout/\fanin}$.

\section{What Manifold Muon Already Enforces}

Manifold Muon is steepest descent under a spectral-norm constraint while the
weights remain on a Stiefel-type constraint set. In its simplest unit-Stiefel
form, it solves for a tangent direction $A$:
\[
    \min_A \langle G, A \rangle
    \quad
    \text{subject to}
    \quad
    \|A\|_2 \le 1,
    \qquad
    A^\top W + W^\top A = 0.
\]
Then it applies
\[
    W^+ = \Retr(W - \eta A).
\]

This construction already controls the same type of quantity used by the
spectral feature-learning condition:
\begin{itemize}
    \item the Stiefel constraint fixes the spectral norm of $W$;
    \item the Manifold Muon tangent problem fixes the spectral norm of $A$;
    \item the tangent constraint removes first-order radial motion;
    \item the retraction returns $W$ to the constraint set after each step.
\end{itemize}

Therefore, with no additional scale factor, unit-Stiefel Manifold Muon gives
\[
    \|W\|_2 = 1,
    \qquad
    \|\Delta W\|_2 \lesssim \eta
\]
up to finite-step and retraction effects.

\section{Why No-Scale Manifold Muon Can Transfer}

The no-scale setting is not arbitrary. It already imposes width-stable
spectral geometry. In square hidden MLP layers, it exactly matches the spectral
feature-learning condition, because
\[
    \fanin = \fanout
    \quad\Longrightarrow\quad
    \sqrt{\frac{\fanout}{\fanin}} = 1.
\]
Thus for a square hidden layer, no-scale Manifold Muon has
\[
    \|W\|_2 = 1,
    \qquad
    \|\Delta W\|_2 \lesssim \eta,
\]
which is precisely the constant-width version of the required spectral scaling.

This gives a direct explanation for good hyperparameter transfer in MLPs whose
width-varying hidden layers are square: the Stiefel constraint and the
unit-spectral tangent constraint already enforce the feature-learning
geometry.

\subsection{Rectangular Input Layers}

The more delicate case is a rectangular layer, for example the first layer of a
CIFAR MLP:
\[
    W \in \R^{H \times 3072}.
\]
If $H < 3072$, unit-Stiefel projection typically gives row-orthonormal weights:
\[
    WW^\top = I_H,
    \qquad
    \|W\|_2 = 1.
\]
The worst-case RMS-to-RMS feature-learning condition would instead ask for
\[
    \|W\|_2
    \asymp
    \sqrt{\frac{H}{3072}}.
\]
So unit-Stiefel no-scale Manifold Muon is not generally identical to the
formal worst-case spectral rule for rectangular layers.

However, row-orthonormal weights have an important average-case property. For
isotropic inputs with unit RMS,
\[
    \mathbb{E}\|Wx\|_2^2 \approx H,
    \qquad
    \rms(Wx) \approx 1.
\]
The same intuition can apply to Manifold Muon tangents. Even if the worst-case
RMS-to-RMS operator norm is too large, the batch RMS of $Ax$ can remain stable
when the input directions are sufficiently spread out.

This suggests the following interpretation:
\begin{itemize}
    \item no-scale Manifold Muon is exactly feature-learning aligned for square
    hidden layers;
    \item rectangular row-Stiefel layers can still preserve typical RMS signal
    and update sizes;
    \item retraction prevents training-time drift away from that geometry.
\end{itemize}

This is a weaker statement than the worst-case spectral condition, but it is a
plausible mechanism behind strong no-scale transfer in current MLP
experiments.

\section{Scaled-Stiefel Feature-Learning Variant}

If the goal is to match the spectral feature-learning condition in worst-case
RMS-to-RMS operator norm, the clean construction is a scaled Stiefel manifold.
Define
\[
    r = \sqrt{\frac{\fanout}{\fanin}},
    \qquad
    W = rQ,
    \qquad
    Q \in \Stiefel.
\]
Then
\[
    \|W\|_2 = r,
    \qquad
    \|W\|_{\rms \to \rms} = 1.
\]

Run the Manifold Muon tangent solve on the unit-Stiefel variable $Q$, producing
a tangent $A$ with $\|A\|_2 \le 1$, and apply the ambient update
\[
    \Delta W = -\eta r A.
\]
Then
\[
    \|\Delta W\|_2 \lesssim \eta r,
    \qquad
    \|\Delta W\|_{\rms \to \rms} \lesssim \eta.
\]
This directly matches the spectral feature-learning requirement for both the
weight and the update:
\[
    \sqrt{\frac{\fanin}{\fanout}}\|W\|_2 = \Theta(1),
    \qquad
    \sqrt{\frac{\fanin}{\fanout}}\|\Delta W\|_2 = \Theta(\eta).
\]

There are two equivalent bookkeeping conventions:
\begin{enumerate}
    \item store the weights on a scaled Stiefel manifold with radius
    $r = \sqrt{\fanout/\fanin}$;
    \item keep unit-Stiefel weights but multiply the tangent update by
    $\sqrt{\fanout/\fanin}$.
\end{enumerate}
The first convention fixes the weight scale. The second fixes the update scale.
When both conventions appear in an implementation, their roles should be made
explicit so the scaling is not accidentally double-counted or omitted.

\section{Relationship to Current Experiments}

The current empirical observation is that no-scale Manifold Muon can show
strong hyperparameter transfer on MLPs. The most conservative explanation is:
\[
    \text{Stiefel weights}
    +
    \text{unit-spectral tangents}
    +
    \text{retraction}
    \quad
    \Rightarrow
    \quad
    \text{stable spectral geometry.}
\]
For square hidden layers, this is the same geometry required by the spectral
feature-learning condition. For rectangular row-Stiefel layers, it may be an
average-case substitute that keeps typical batch RMS changes stable.

This leads to two separate claims:
\begin{description}
    \item[No-scale MLP Manifold Muon.]
    A simple geometry-stabilized optimizer. It can transfer well because square
    hidden layers already satisfy the spectral rule and rectangular
    row-Stiefel layers can preserve typical RMS behavior.

    \item[Feature-learning MLP Manifold Muon.]
    A scaled-Stiefel or RMS-to-RMS-normalized variant. This is the clean
    worst-case operator-norm construction aligned with the spectral
    feature-learning theory.
\end{description}

\section{What to Measure Next}

To distinguish the exact spectral story from the average-case Stiefel story,
measure the following per layer during width and aspect-ratio sweeps:
\begin{itemize}
    \item $\|W\|_2$;
    \item $\|\Delta W\|_2$;
    \item $\sqrt{\fanin/\fanout}\|W\|_2$;
    \item $\sqrt{\fanin/\fanout}\|\Delta W\|_2$;
    \item $\rms(Wx) / \rms(x)$ on real minibatches;
    \item $\rms(\Delta W x) / \rms(x)$ on real minibatches;
    \item tangent residual $\|W^\top A + A^\top W\|$;
    \item Stiefel residual after retraction.
\end{itemize}

The decisive test is to vary aspect ratio, not only width. A useful MLP sweep
should include:
\begin{itemize}
    \item square hidden layers $H \times H$;
    \item expanding layers $4H \times H$;
    \item contracting layers $H \times 4H$;
    \item fixed-input layers $H \times d_{\mathrm{in}}$.
\end{itemize}

If no-scale Manifold Muon transfers only in square or row-Stiefel regimes, the
mechanism is likely spectral/Stiefel geometry plus average-case RMS
preservation. If it transfers across changing aspect ratios while the
RMS-to-RMS operator norm is unstable, then measured batch geometry may be more
relevant than worst-case operator norm for these MLP experiments.

\section{References}

\begin{itemize}
    \item \href{https://arxiv.org/abs/2310.17813}{A Spectral Condition for Feature Learning}.
    \item \href{https://arxiv.org/abs/2405.14813}{Scalable Optimization in the Modular Norm}.
    \item \href{https://thinkingmachines.ai/blog/modular-manifolds/}{Manifold Muon derivation}.
\end{itemize}

\end{document}
